\section{Related Work}
\label{sec:related_work}

\textbf{Parameter-Efficient Fine-Tuning \& LoRA Serving.} Low-Rank Adaptation (LoRA) \cite{hu2021lora} has emerged as the dominant paradigm for adapting large pre-trained models to specific tasks by parameterizing weight updates via low-rank matrices. To deploy hundreds of fine-tuned LoRA models on a single GPU node, multi-tenant serving frameworks like S-LoRA \cite{sheng2023slora} and Punica \cite{chen2023punica} batch request sequences and swap adapter weights dynamically. While these frameworks optimize hardware execution, they typically serve separate tasks on isolated request streams. In contrast, test-time dynamic model merging focuses on multi-task scenarios where individual queries in a single stream require joint or mixed expert contributions.

\textbf{Static and Dynamic Model Merging.} Model merging combines multiple fine-tuned models into a single set of parameters to preserve multi-task capabilities without increasing inference cost. Static methods like Model Soups \cite{wortsman2022model}, Task Arithmetic \cite{ilharco2022editing}, and Ties-Merging \cite{yadav2023ties} compute a fixed, offline average or projection of parameter weights. To handle complex non-linear representation shifts at test-time, dynamic model ensembling has moved toward sample-wise activation blending. SABLE \cite{sable_2024} and SPS-ZCA \cite{zhang24spszca} use representation signals extracted from early layers to determine task coordinates and blend LoRA activations in subsequent layers in a single pass. However, these methods are stateless: they treat sequential queries independently, leading to extreme routing jitter and representation instability.

\textbf{Stateful Serving \& Control Theory.} To mitigate routing jitter, stateful ensembling has recently emerged as a design paradigm. ChemMerge \cite{chemmerge_2026} treats model ensembling as a continuous-time multi-component chemical reactor, using first-order reaction kinetics ordinary differential equations (ODEs) to smooth routing trajectories over sequential inputs. While ChemMerge empirically reduces routing weight jitter, it is a heuristic control strategy without any learning-theoretic or stability guarantees. Its parameters are not learned but rather statically set, causing severe accuracy drops under heterogeneous workloads. PAC-Kinetics bridges this gap by unifying stateful kinetics with contractive, input-to-state stable control theory and learning-theoretic optimization.

\textbf{PAC-Bayesian Generalization Theory.} PAC-Bayesian theory, originally introduced by McAllester \cite{mcallester1999pac, mcallester2003pac} and extensively developed by Catoni \cite{catoni2007pac} and others \cite{guedj2019primer, alquier2021user}, provides rigorous out-of-sample generalization guarantees for randomized predictors. Standard PAC-Bayesian bounds rely heavily on the assumption that training and test samples are independent and identically distributed (i.e.d.). In streaming and dynamic edge serving, this assumption is frequently violated due to temporal correlations and non-stationary task shifts. To address this, we build upon the theoretical foundations of PAC-Bayesian bounds for stationary mixing processes \cite{alquier2013pac}, establishing the first mathematically rigorous generalization bound for stateful, sequential model ensembling.
